\section{Rollen, Verantwortlichkeiten und Governance}

In der Praxis stellt sich die Frage, wer Risiken erkennt, bewertet, freigibt und nachsteuert.
Governance liefert dafür den organisatorischen Rahmen: Ziele und Entscheidungen müssen fachlich begründet sein,
Daten und Modelle kontrolliert erstellt und betrieben werden, regulatorische und ethische Anforderungen laufend eingehalten.

Im DASC-PM ist Governance ein interdisziplinäres Zusammenspiel klar abgegrenzter Rollen und Entscheidungspunkte \cite{schulz_dasc-pm_2022}.
Diese Rollentrennung reduziert methodische Fehlinterpretationen, Compliance-Verstöße und wirtschaftliche Fehlinvestitionen \cite{Saltz2015}.

\subsection{Rollendefinition und interdisziplinäre Verantwortlichkeit}
Das DASC-PM geht nicht davon aus, dass eine einzelne Rolle alle fachlichen, methodischen und technischen Aufgaben gleichzeitig trägt \cite{Saltz2015,schulz_dasc-pm_2022}.

Stattdessen werden Verantwortlichkeiten funktional aufgeteilt und Schnittstellen klar beschrieben.

Im Kern lassen sich drei Verantwortungsbereiche unterscheiden:
\begin{itemize}
  \item \textbf{Fachliche Verantwortung durch Domänenexpert:in oder Business Owner:}
        Ziel- und Nutzenorientierung, Validierung des Business Value und Entscheidung über den Einsatz im Geschäftsprozess \cite{schulz_dasc-pm_2022}.
  \item \textbf{Methodische Verantwortung durch Data Scientist:}
        Auswahl geeigneter Verfahren, Evaluation/Dokumentation der Modellgüte und Minimierung methodischer Risiken \cite{schulz_dasc-pm_2022}.
  \item \textbf{Technische Verantwortung durch Data Engineer, IT oder Plattformbetrieb:}
        Datenintegrität, Nachvollziehbarkeit der Datenflüsse, Betriebsstabilität und Sicherheit entlang von Data- und IT-Governance \cite{Sculley2015}.
\end{itemize}

Zur klaren Zuweisung von Verantwortlichkeiten kann eine RACI-Matrix genutzt werden; COBIT empfiehlt solche Zuordnungen als Bestandteil wirksamer IT-Governance \cite{ISACA2018}.

\begin{table}[htbp]
\centering
\small
\setlength{\tabcolsep}{4pt}
\renewcommand{\arraystretch}{1.1}
\caption{Vereinfachtes Beispiel einer RACI-Zuordnung für zentrale Governance-Entscheidungen
in Data-Science-Projekten.
A steht für Accountable, R für Responsible, C für Consulted, I für Informed.
Eigene Darstellung nach \cite{ISACA2018}}
\label{tab:raci-ds}
\begin{tabularx}{\textwidth}{@{}>{\raggedright\arraybackslash}Xcccc@{}}
\toprule
\textbf{Governance-Objekt}
  & \textbf{\shortstack{Domäne}}
  & \textbf{\shortstack{Data\\Scientist}}
  & \textbf{\shortstack{DE/\\IT}}
  & \textbf{\shortstack{Compliance\\/DSB}} \\
\midrule
Problem- \& Zieldefinition      & A/R & C & I & C \\
Datenzugang \& Datenqualität      & A   & C & R & C \\
Modellierung \& Evaluation         & A   & R & C & C \\
Deployment \& Betrieb            & A   & C & R & C \\
Go/No-Go \& Nutzung im Prozess & A   & C & C & R \\
\bottomrule
\end{tabularx}
\end{table}

\subsection{Governance als Risikomanagement-Struktur}
Governance unterstützt das Risikomanagement besonders dann, wenn sie neben Rollen auch Kontroll- und Entscheidungspunkte entlang des Projektverlaufs etabliert.

Diese Logik passt zur Iteration des DASC-PM: Rücksprünge und formale Entscheidungspunkte verhindern, dass Projekte trotz fehlender Eignung oder fehlendem Nutzen weiterlaufen \cite{schulz_dasc-pm_2022}.

Drei Governance-Ebenen sind hierfür besonders relevant:
\begin{itemize}
  \item \textbf{Entscheidungs-Governance:}
    Management- oder Business-Entscheidungen an kritischen Punkten halten das Projekt auf Nutzen und Strategie ausgerichtet \cite{schulz_dasc-pm_2022}.
  \item \textbf{Compliance-Governance im Sinne von ''by Design,,:}
    Regulatorische Anforderungen wie Datenschutz, Informationssicherheit und KI-Regulierung müssen von Beginn an mitlaufen \cite{EUAIAct2024}.
  \item \textbf{Ethische Governance:}
    Fairness, Transparenz, Erklärbarkeit und die Vermeidung diskriminierender Effekte sind als Governance-Aufgaben fest zu verankern \cite{FloridiCowls2019}.
\end{itemize}

\subsection{Geteilte Verantwortung, Accountability und Haftung}
Ein zentrales Governance-Problem datengetriebener Systeme ist die Verantwortung bei Fehlentscheidungen: methodische Korrektheit und organisatorische Angemessenheit fallen häufig auseinander.

Das DASC-PM löst diese Spannung über geteilte Verantwortung: Der Data Scientist verantwortet die methodische Qualität, die fachliche Instanz die Anwendung, Grenzen und Interpretation im Geschäftsprozess \cite{schulz_dasc-pm_2022}.

Trotz geteilten Rollen muss die Rechenschaftspflicht klar sein. Accountability bedeutet, dass Organisationen Verantwortliche und Nachweisbarkeit über Entscheidungen, Datenstände und Modellversionen festlegen \cite{FloridiCowls2019}.

COBIT trennt hierfür Governance (evaluate, direct and monitor) und Management (plan, build, run and monitor) als Leitlinie für die Trennung von Entscheidungsverantwortung und operativer Umsetzung \cite{ISACA2018}.
